\begin{abstract}
This research studies whether the predictability of next-month U.S. common-stock returns varies across high- and low-volatility market regimes and whether nonlinear machine-learning methods capture this variation differently from a regularized linear benchmark. The empirical pipeline uses monthly CRSP stock data, a VIX-based regime classification constructed from an expanding historical median, and strictly time-respecting out-of-sample estimation with Elastic Net and XGBoost. The current preliminary specification uses a computationally reduced fixed large-cap universe and a partial predictor set consisting of size and momentum; book-to-market is added in the subsequent Compustat/CCM data-construction stage. Over 83 out-of-sample formation months from January 2019 through November 2025, the tuned Elastic Net collapses to an intercept-only cross-sectional forecast and therefore provides no usable stock-ranking signal. XGBoost retains cross-sectional forecast dispersion, but its mean monthly Spearman rank correlation with realized next-month returns is slightly negative. Equal-weighted XGBoost D10-minus-D1 returns are close to zero overall, whereas the value-weighted spread averages approximately 0.71\% per month and is concentrated in low-VIX months, where it averages approximately 1.43\% per month; the corresponding high-VIX value-weighted spread is essentially zero. These findings are descriptive and preliminary rather than formal evidence of statistical or net-of-cost profitability. They motivate the next stages of the research, including the complete Compustat/CCM predictor set, formal inference, turnover measurement, and transaction-cost robustness.
\end{abstract}
